% Independent supporting proof fragment. No standalone document.
\section{A controlled dispersion stress under the unchanged GBM class}
\label{supp:dispersion-scaling-proof}

I vary a financial exposure parameter while retaining the model class,
common Gaussian smoothing scale, forward carry, contract and numerical
law. This is a parameter stress of a classical risk-neutral model. It is
not a new stochastic model or a new no-arbitrage principle.

\subsection{Exact model and ratio invariants}

Let $\theta=r_f-q$, let $R_i\in\mathbb R^p$ be the original outer
exposure of asset $i$, and fix the strictly positive common exposure
$\lambda$. The frozen maturity is one. For $c>0$, the stressed terminal
asset is
\[
 S_i^{(c)}=K\exp\!\left(\mu_i(c)+cR_i\cdot Y+\lambda Z\right),
 \qquad
 \mu_i(c)=\theta-\frac12\{c^2\|R_i\|^2+\lambda^2\},
\]
where $Y$ has independent standard Gaussian coordinates and $Z$ is an
independent standard Gaussian. All original exposures, weights and
the selected $c\in\{1,3/2,2,3\}$ are exact rational inputs.
Consequently the recomputed drift and covariance entries are exact
rational quantities. For positive basket weights of sum one, put
\[
 R_G=\sum_iw_iR_i,\quad Q_R=\sum_iw_i\|R_i\|^2,\quad
 V_R=Q_R-\|R_G\|^2=\sum_iw_i\|R_i-R_G\|^2\geq0.
\]
Define $m_G(c)=\sum_iw_i\mu_i(c)$,
$d_i(c)=\mu_i(c)-m_G(c)$ and
$s_i(c)=cR_i-cR_G$. Then exact subtraction gives
\begin{align*}
 m_G(c)&=\theta-\tfrac12\lambda^2-\tfrac12c^2Q_R
       =m_G(1)-\tfrac12(c^2-1)Q_R,\\
 d_i(c)&=-\tfrac12c^2(\|R_i\|^2-Q_R)=c^2d_i(1),\\
 s_i(c)&=cs_i(1),\qquad R_G(c)=cR_G(1).
\end{align*}
The common scale and forward carry cancel from the drift differences;
neither is multiplied by $c$. Weighted centring remains exact:
$\sum_iw_id_i(c)=0$ and $\sum_iw_is_i(c)=0$.

The terminal covariance matrix is
$\Sigma(c)=c^2RR^{\mathsf T}+\lambda^2\mathbf1\mathbf1^{\mathsf T}$.
For every $x\in\mathbb R^n$,
\[
 x^{\mathsf T}\Sigma(c)x
 =c^2\|R^{\mathsf T}x\|^2+\lambda^2(\mathbf1^{\mathsf T}x)^2\geq0.
\]
Thus this covariance is valid without requiring nonsingularity. Gaussian
integration proves
$\mathbb E_Q S_i^{(c)}=K e^\theta$ for every asset and stress scale.
The arithmetic forward $\mathbb E_Q A^{(c)}=K e^\theta$ is unchanged.
The geometric forward is
\[
 \mathbb E_Q G^{(c)}=K\exp\{\theta-c^2V_R/2\},\qquad
 \operatorname{Var}\log G^{(c)}=c^2\|R_G\|^2+\lambda^2>0.
\]
These are separate expectations. Dividing them does not generally give
$\mathbb E_Q(A^{(c)}/G^{(c)})$, because the ratio and geometric basket
are dependent.

For completeness, the terminal law is generated by the standard
constant-coefficient dynamics
\[
 \frac{dS_{i,t}^{(c)}}{S_{i,t}^{(c)}}=
 (r_f-q)dt+c\sum_jR_{ij}dB_{j,t}+\lambda dB_{0,t},
 \qquad S_{i,0}^{(c)}=K,
\]
under the stipulated pricing measure $Q$. The independent Brownian
coordinates integrate to the displayed terminal law. The exponential
martingale has expectation one because its deterministic quadratic
variation is finite. It follows that $e^{-(r_f-q)t}S_{i,t}^{(c)}$ is
a true martingale. Equivalently, the discounted gain process
$e^{-r_ft}S_{i,t}^{(c)}+\int_0^t e^{-r_fs}qS_{i,s}^{(c)}ds$ is a
martingale. Under the project's stated equivalent-pricing-measure
assumption, this supplies the usual no-arbitrage pricing condition for
admissible self-financing strategies. The stress does not prove
equivalence to an unspecified physical measure or introduce a new
empirical risk premium.

Let $X_i(c,y)=\exp\{c^2d_i(1)+cs_i(1)\cdot y\}$ and
$T_c(y)=\sum_iw_iX_i(c,y)$. Jensen gives
\[
 T_c(y)\geq\exp\!\left(\sum_iw_i\{c^2d_i(1)+cs_i(1)\cdot y\}\right)=1.
\]
Hence $U_c=T_c-1\geq0$ for every real $y$, at every selected scale.
The arithmetic basket remains $A^{(c)}=G^{(c)}T_c(Y)$, and common
Gaussian conditioning gives the same scalar kernel $P_\lambda$.
The residual log path $1+zU_c$ is at least one on $z\in[0,1]$.
Thus the existing residual derivative envelopes apply after their exact
geometry and geometric loading are recomputed. Their coefficients
$M_k,D_1,D_2,D_{31}$ remain the same payoff-specific conditional
norms at the unchanged $\lambda$. Their geometric moments do not remain
the same and cannot be borrowed from the original exposure.

\subsection{An unchanged retained Gaussian law}

The latent variables remain standard Gaussian. At unchanged
$L=6,N=2^{10},q=28$, the continuous retained measure $\nu$, actual
rounded measure $\tau$ and all one-coordinate transport measures are
unchanged. Their mass $a=\Phi(L)-\Phi(-L)$ and the product retained
mass $a^p$ depend on the latent dimension, but not on $c$ or the financial
drift. Every hybrid moment retains the raw factors of $a$ through its
scalar transforms. There is no scale-dependent renormalisation of the
Gaussian label distribution.

For a signed geometry polynomial
$h_c(y)=\sum_\alpha b_\alpha X(c,y)^\alpha$, its rate and constant
are respectively $c\sum_i\alpha_is_i(1)$ and
$c^2\sum_i\alpha_id_i(1)$. Full Gaussian moments therefore equal
\[
 \mathbb E h_c(Y)=\sum_\alpha b_\alpha
 \exp\!\left[c^2\left\{\sum_i\alpha_id_i(1)
                    +\frac12\left\|\sum_i\alpha_is_i(1)\right\|^2\right\}\right].
\]
The hybrid moment uses the same rates in the unchanged positive
transport and continuous/rounded scalar transforms. The Gaussian
outside-box monomial moment is
\[
 e^{c^2(b+\|v\|^2/2)}
 \left[1-\prod_j\{\Phi(L-cv_j)-\Phi(-L-cv_j)\}\right],
 \quad b=\sum_i\alpha_id_i(1),\quad v=\sum_i\alpha_is_i(1).
\]
Both formulas retain signed polynomial coefficients; positivity is
established from the original $U_c$ geometry, not by replacing those
coefficients with absolute values. The outer tail is charged in full.

For $k\geq1$, the price coefficient is
$c_k(c)=\mathbb E[U_c^k C_{k,\lambda}(m_G(c)+cR_G\cdot Y)]/k!$,
with $C_{1,\lambda}=P_\lambda'$. Completing the Gaussian square in a
monomial term gives the exact shifted convolved kernel at scale
$\sigma(c)=\sqrt{\lambda^2+c^2\|R_G\|^2}$ and mean
$m_G(c)+c^2R_G\cdot\sum_i\alpha_is_i(1)$, multiplied by its full
Gaussian exponential moment. The drift shift and covariance shift both
have order $c^2$. Omitting the latter would alter the coefficient whenever
$R_G\ne0$.

\subsection{Why a fixed Taylor degree can lose usefulness}

Suppose $V_R>0$. For integer $k\geq1$, Gaussian integration and two
applications of Jensen give
\begin{align*}
 \mathbb E T_c
 &=\sum_iw_i\exp\{c^2(d_i(1)+\|s_i(1)\|^2/2)\}
   \geq e^{c^2V_R/2},\\
 \mathbb E U_c^k&\geq(\mathbb E U_c)^k
   \geq\{e^{c^2V_R/2}-1\}^k.
\end{align*}
The first inequality uses
$\sum_iw_id_i(1)=0$ and
$\sum_iw_i\|s_i(1)\|^2=V_R$; the second uses $U_c\geq0$.
For any fixed degree $m$ and positive fixed norm $M_{m+1}$, the raw
global finite-Taylor radius obeys
\[
 \frac{M_{m+1}}{(m+1)!}\mathbb E U_c^{m+1}
 \geq\frac{M_{m+1}}{(m+1)!}
          \{e^{c^2V_R/2}-1\}^{m+1}\longrightarrow\infty.
\]
This is an explicit limitation of the fixed conditional-norm radius.
Intersecting a capped interval with $[0,H]$ prevents its final width
from diverging, but need not restore the requested small monetary
radius. The result does not show that the true price becomes difficult,
that every classical method fails, or that the transport allowance stays
small. Its own higher dispersion moments and tail can also grow.
The divergence does not apply to an intersection of independent valid
intervals or to a different payoff-specific remainder. Indeed,
$0\leq\mathbb E(A^{(c)}-K)_+\leq K e^\theta$ at every scale,
and $0\leq\mathbb E F_H(A^{(c)})\leq H$. The underlying price ranges
remain bounded while the displayed raw Taylor radius diverges.

For every fixed finite stress $c$, the existence of arbitrarily accurate
mathematical transfer certificates is an established consistency fact.
The lognormal ratio has a finite second moment. Thus dominated
convergence gives
$\mathbb E[U_c^2\mathbf1_{Y\notin[-L,L]^p}]\to0$ as $L\to\infty$.
Fixing $L$ afterwards, let $\delta$ be the maximum cell width.
Equal-CDF cells satisfy
$\delta\leq a/(N\min_{[-L,L]}\phi)\to0$.
For any actual node in its cell,
\[
 \rho_1(\mathbb R)\leq a\delta,\qquad
 \rho_2(\mathbb R)\leq a\delta^2/2,
\]
where the first assertion uses the exact dominating atom weights
$|\mu_{1,c}|$. Arbitrarily tight rational upper weights add an
arbitrarily small numerical allowance. Every residual derivative and
geometry envelope is bounded on this compact box. The hybrid
allowances, including their homogeneous positive-measure Cauchy
closures, therefore tend to zero with these transport masses.
For midpoint quantiles the distance to either cell endpoint is at least
$a/(2N\max\phi)$. A dyadic approximation error smaller than that
distance preserves node containment; $q$ can be increased to satisfy
this requirement. Any rounding tie still has two admissible nearest
nodes at sufficiently fine precision, although an actual implementation
must certify its chosen rounding convention. Outward precision can be
increased to certify the remaining numerical allowances.
Choosing $L$ for the tail and then $N,q$ for the compact transfer proves
consistency. It is not a claim that the frozen $L=6,N=1024,q=28$ passes
at every stress, or that the required bit counts and physical resources
are small. In particular, failure of a fixed-degree raw Taylor radius
does not imply absence of a transfer certificate.

I compare the stated transfer allocation $0.005$ with a complete
classical price screen of radius $0.01$. These certify different
objects. A transfer allowance does not determine a price until the
remaining residual, arithmetic and estimator budgets have been met.
A case admitted by transport but rejected by a degree-three or
degree-five price screen identifies a limitation of that specified
screen. It is neither an algorithmic separation theorem nor evidence
of physical quantum advantage. Higher-degree checks introduce new
degree-five/six moments and payoff-kernel evaluations, whose costs and
information must be recorded. All original scale-one results and every
stressed failure must remain visible.
